Modern life feels unhinged. Modern technologies drive me crazy; every day I am flooded with loud, trivial, and overwhelming news. The United States edges toward war with Iran. Energy is becoming more expensive while Germany, having shut down its nuclear power stations years ago, now finds it difficult to restart them. France struggles to form a stable government, even as its deficits need to be paid.

I have no desire to escape to an Asian monastery in search of peace. A life that simple might be quiet, but for me it would be a waste of time. Instead, I want to focus on something that does not change, something that, ideally, only I can do. The human brain grasps stories easily and tries to make sense of them, even when there is no real sense to be found. At the same time, the human brain is notoriously bad at statistics and probabilities.

Casanova’s life is a useful example. It is tempting to look only at his adventures, but what interests me is that he helped the French government sell state lotteries to raise money, and he made a fortune doing it. Later, he lost that fortune at the gambling tables. Both his rise and his ruin were tied to probability.

Ed Thorp managed to beat the blackjack dealers using statistics, at a time when his peers insisted it was impossible. Even scientists—people who should have known better—were making confident, unscientific claims.

Now I live in Greece, where people play the lottery constantly. Everyone repeats that it is impossible to beat, and their certainty reminds me of Thorp’s colleagues. Yet at the heart of every lottery sits a random number generator, and designing a truly good one is much harder than most people think.

Knowledge of statistics and probability is always useful. I hope I can use it not only to understand these games more clearly, but also, perhaps, to earn some money.

\vspace{2em}
\begin{flushright}
Ming Lu\\
\emph{Aegina, Greece}\\
\emph{December 19, 2025}
\end{flushright}

